% ============================================================================
%  SECCIÓN 1.11: MÉTODOS DE INVESTIGACIÓN
% ============================================================================

\section{M\'etodos de investigaci\'on}

Los m\'etodos de investigaci\'on se clasifican en tres grandes grupos, de los
cuales el investigador selecciona los que resulten m\'as convenientes para su
investigaci\'on: m\'etodos te\'oricos, m\'etodos emp\'iricos y m\'etodos
estad\'isticos \citep{arias2012, hernandez2014}.

\begin{table}[H]
  \centering
  \caption{M\'etodos de investigaci\'on considerados en el proyecto}
  \label{tab:metodos}
  \contenidoConFuente{%
    \begin{tablaAPA}
      \begin{tabular}{@{}p{3.4cm}p{5.6cm}p{6.8cm}@{}}
        \toprule
        \textbf{Clase de m\'etodo} & \textbf{M\'etodos} &
        \textbf{Uso en el proyecto} \\
        \midrule
        Te\'oricos &
        Hist\'orico l\'ogico; an\'alisis y s\'intesis; inducci\'on y deducci\'on;
        modelaci\'on &
        Revisi\'on de antecedentes, an\'alisis de la bibliograf\'ia y construcci\'on
        del marco te\'orico \\
        \addlinespace
        Emp\'iricos &
        Observaci\'on; entrevista; encuesta; revisi\'on documental &
        Recolecci\'on de informaci\'on sobre los escenarios deportivos y las
        necesidades de los usuarios \\
        \addlinespace
        Estad\'isticos &
        Estad\'istica descriptiva &
        Procesamiento y an\'alisis de los resultados de las encuestas aplicadas
        al p\'ublico \\
        \bottomrule
      \end{tabular}
    \end{tablaAPA}%
  }{Elaboraci\'on propia en base a la revisi\'on documental.}
\end{table}

\subsection{M\'etodo hist\'orico l\'ogico}

Este m\'etodo se utilizar\'a para estudiar la trayectoria de la infraestructura
deportiva del pa\'is y la evoluci\'on de los medios de difusi\'on de la oferta
deportiva, partiendo de los antecedentes hist\'oricos para llegar a la
comprensi\'on de la situaci\'on actual del objeto de estudio.

\subsection{M\'etodo anal\'itico--sint\'etico}

Se emplear\'a para descomponer la problem\'atica en sus elementos (causas y
efectos), analizar cada uno de ellos y sintetizar las conclusiones que
fundamenten la propuesta de soluci\'on.

\subsection{Encuesta y entrevista}

La encuesta se aplicar\'a al p\'ublico general para conocer sus h\'abitos
deportivos, el acceso a la informaci\'on y la aceptaci\'on de la soluci\'on
propuesta. La entrevista se realizar\'a a los gestores de los escenarios
deportivos para identificar sus necesidades de visibilidad y los requisitos
del sistema desde la perspectiva institucional \citep{hernandez2014}.
